% -------- %
% Packages %
% -------- %

%% Basics
\usepackage[french]{babel} % langue
\usepackage[utf8]{inputenc} % means input encoding. A voir avec l'encodage des caractères. faut pas chercher à comprendre.
\usepackage[T1]{fontenc} % better accents, and unicode caracters
\usepackage{lmodern} % pareil

%% figures
\usepackage{graphicx} % Required for inserting images
\usepackage{caption}
\usepackage{subcaption} % Required for creating figures with multiple parts (subfigures)
\usepackage{float} % Allows putting an [H] in \begin{figure} to specify the exact location of the figure
\usepackage{wrapfig} % Allows in-line images such as the example fish picture
\usepackage{adjustbox} % Required for fitting table to page
\usepackage{mathtools} % Arrows with text above and below
\graphicspath{{./images/}}		%chemin vers les images

%% tables
\usepackage[table]{xcolor} % Required for specifying colors by name. table allows to color in tables.
\usepackage{booktabs} % Required for better table rules : \toprule, \midrule, \bottomrule and set the width of the columns
\usepackage{multirow} % Required for multirows in tables
\usepackage{tikz} % Required for drawing custom shapes

%% Math
\usepackage{amsmath} % Required for some math elements
\usepackage{amsthm} % to write theorems, definitions, lemmas, ...
\usepackage{amssymb} % more math symbols

%% mise en page
\usepackage{geometry} % Required for adjusting page dimensions and margins. Adapted the margin for comments
\usepackage{hyperref} % manage hyperlinks
\usepackage{xspace} % space if no punctiation
\usepackage{xcolor} % Required for specifying colors by name
\definecolor{darkgreen}{rgb}{0.0, 0.5, 0.0}
\geometry{margin=20mm} %enough space for margin notes
\hypersetup{colorlinks=true, linkcolor=darkgreen, urlcolor=magenta, citecolor=blue}

% % ------------------- %
% % Writing pseudocodes %
% % ------------------- %
\usepackage{algorithm} % write pseudocode, suggested by IEEEtran instead of algorithm2e
\usepackage{algpseudocode} 
% \usepackage{algorithmic}

%% TODO notes
\usepackage{todonotes} % Required for the boxes that can be used to write comments, and to-do lists

%% Subsubsubsection
\usepackage{titlesec}

\titleclass{\subsubsubsection}{straight}[\subsection]

\newcounter{subsubsubsection}[subsubsection]
\renewcommand\thesubsubsubsection{\thesubsubsection.\arabic{subsubsubsection}}
\titleformat{\subsubsubsection}
  {\normalfont\normalsize\bfseries}{\thesubsubsubsection}{1em}{}
\titlespacing*{\subsubsubsection}
{0pt}{3.25ex plus 1ex minus .2ex}{1.5ex plus .2ex}

\makeatletter
\def\toclevel@subsubsubsection{4}
\def\l@subsubsubsection{\@dottedtocline{4}{7.0em}{4.1em}}
\makeatother

\setcounter{secnumdepth}{4}
\setcounter{tocdepth}{4}

\newcommand{\weightedlinkright}[1]{$\xrightarrow{\text{\hspace{0.2cm}\scriptsize{#1}\hspace{0.2cm}}}$} % Intermediate state, should not be used directly
\newcommand{\link}[3]{\mbox{#1 \weightedlinkright{#2} #3}} % link from #1 to #3 with weight #2

\usepackage{etoolbox}

\newcounter{timecounter}

\newcommand{\evolvinglink}[3]{
    \setcounter{timecounter}{0}
    \renewcommand*{\do}[1]{%
        \begin{minipage}[b]{0.22\textwidth}
            \centering
            \begin{minipage}{\textwidth}
                \centering
                \begin{tikzpicture}[shorten >=1pt, auto, node distance=18mm]
                    \begin{scope}[every node/.style={circle,draw=black}]
                        \node (0) {0};
                        \node (1) [right of=0] {1};
                    \end{scope}
                    \begin{scope}[every edge/.style={draw=black,ultra thick}]
                        \ifnumcomp{##1}{=}{-1}
                            {} % Do nothing if the weight is -1
                            {\draw[->] (0) -- node[above] {##1} (1);} % Draw the edge otherwise
                    \end{scope}
                \end{tikzpicture}
            \caption*{$#1_{\thetimecounter}$}
            \end{minipage}
        \end{minipage}
        \stepcounter{timecounter}
    }
    \docsvlist{#2} % Process list of weights
    \caption*{#3}
}

\usepackage{etoolbox}

\newcommand{\indexinge}[2]{
    %index into the list of nodes
    %\foreach \elem [count=\i] in #1 {
        %\ifnum\i=#2
            %\elem
        %\fi
    %}

    % Create a new empty list
    \newcommand{\nodes}{}
    % Fill the list with the elements of the first argument
    \foreach \elem [count=\i] in #1 {
        % if i is equal to the second argument, add the element to the list
        \ifnum\i=#2
            \xappto{\nodes}{\elem}
        \fi
    }

    % Print the list
    \nodes

}

\newcommand{\indexing}[2]{
    #1[#2]
}

\newcommand{\indexingarray}[2]{
    #1(#2)
}


\newcommand{\graphpath}[2]{
        \def\nodes{#1}
        \def\weights{#2}
        \foreach \node [count=\i] in \nodes {
            \ifnum\i=1
                \node
                \weightedlinkright{\weights[0]}
            \else
                \link{\node}{\weights[\i-2]}{\nodes[\i]}
            \fi
        }
}

\newcommand{\graphpatheee}[1]{
    \def\nodes{#1}
    \foreach \node [count=\i] in \nodes {
        % if i is pair (modulo 2)
        \ifodd\i
        \else
            \weightedlinkright{\node}
        \fi
    }
}

%\usepackage{luacode}

%\newcommand{\indexeee}[2]{
    %\directlua{
        %local list = {#1}
        %tex.sprint(list[#2])
    %}
%}